% Source: https://artofproblemsolving.com/community/c3h1319011_math_courts_warmup_3_26
% Stable identity: aops-topic-1319011-q26
The scale factor from $ABC$ to $XYZ$ is the ratio of hypotenuses, $\dfrac{18}{10}=\dfrac{9}{5}$. Areas scale by the square of the scale factor, so the ratio is $\left(\dfrac{9}{5}\right)^2=\boxed{\dfrac{81}{25}}$. (Directly: $[ABC]=24$ and $XYZ$ has legs $10.8$ and $14.4$, area $77.76=\dfrac{81}{25}\cdot24$.)

\begin{asy}
usepackage("amsmath");
size(6cm);
defaultpen(fontsize(10pt));
path ra(pair P,pair u,pair v,real s=0.2){return P+s*unit(u)--P+s*unit(u)+s*unit(v)--P+s*unit(v);}
pair A=(0,0),B=(4,0),C=(0,3); draw(A--B--C--cycle,linewidth(0.9)); draw(ra(A,B-A,C-A,0.25));
label("$6$",(A+C)/2,W); label("$8$",(A+B)/2,S); label("$10$",(B+C)/2,NE);
pair X=(6,0),Y=(6+7.2,0),Z=(6,5.4); draw(X--Y--Z--cycle,linewidth(0.9)); draw(ra(X,Y-X,Z-X,0.25));
label("$18$",(Y+Z)/2,NE); label("$ABC$",(1.2,0.9)); label("$XYZ$",(8.2,1.5));
label("\mbox{}",(min(currentpicture,user=true)+max(currentpicture,user=true))/2,black); // keeps the text after the figure black when LaTeX typesets the labels inline
\end{asy}
